\documentclass[12pt]{article}
\usepackage[utf8]{inputenc}
\usepackage[T1]{fontenc}
\usepackage{hyperref}
\title{Garrwa Language Resources for K-12 Curriculum Development}
\author{Compiled by Claude Sonnet 5.5 (claude-sonnet-5-5) for Andr\'as Kornai}
\date{}
\begin{document}
\maketitle

\section{Introduction}
Garrwa (ISO 639-3 wrk; Glottocode gara1269; Garrwan family) is spoken around Borroloola and Robinson River in the Gulf of Carpentaria, Northern Territory. It is severely endangered. The 2006 census counted about 35 people in the Borroloola area who spoke it at home. Very little curriculum-ready material turned up in search. What exists is mainly a recent plants-and-animals book, a reference grammar and some archival primers.

\section{Primary Educational Resources}
\begin{itemize}
\item \emph{Garrwa and Gunindirri Garrwa Plants and Animals} (2021), by the Garrwa Language Team with Ilana Mushin and Glenn Wightman. It records 253 plants and 279 animals. Most entries have a QR code linking to audio by Daphne Mawson. A UQ research listing gives a slightly different title (``Garrwa, Gunindirri and Gudanji Plants and Animals''). The news item is at \url{https://languages-cultures.uq.edu.au/node/8953}.
\item Archival Garrwa (Karawa) primers and bilingual texts in the Living Archive of Aboriginal Languages, held on Territory Stories. Search results mention a 1982 trial primer and a 1995 bilingual text, \emph{Camping in Caves}. The Garrwa-language items are at \url{https://territorystories.nt.gov.au/10070/871570}. I could not confirm which record is which.
\end{itemize}

\section{Linguistic and Reference Materials}
\begin{itemize}
\item Ilana Mushin, \emph{A grammar of (Western) Garrwa} (De Gruyter Mouton, 2012). I saw no open-access URL for it in the results, and it is not listed in the json.
\item The AIATSIS AUSTLANG entry for Garrwa is at \url{https://collection.aiatsis.gov.au/austlang/language/c23}.
\item The Endangered Languages Project entry is at \url{https://endangeredlanguages.com/lang/4190}.
\item The Wikipedia article is at \url{https://en.wikipedia.org/wiki/Garrwa_language}.
\end{itemize}

\section{Religious and Community Resources}
I found no scripture or community-organisation pages in the search results. The Garrwa Language Team is named as the author of the 2021 book. I saw no contact details and give none.

\section{Digital and Audio Resources}
Audio exists as QR-linked recordings of Daphne Mawson in the 2021 book. I saw no standalone URL for the audio.

\section{Curriculum Development Context}
Garrwa is not known to be taught as a regular school subject. The 2021 book and the Living Archive primers are the most usable starting points. The search surfaced no dictionary, so curriculum work would need to build one with the Garrwa Language Team and elders. Any use of archival items should respect their CC BY-NC-ND 4.0 licence. Remaining material is thin, and a dossier cannot honestly claim more.
\end{document}